\chapter{引言}


\section{研究背景}

新课改背景下，教育部颁布了《普通高中数学课程标准（2017年版）》。课程标准第一次提出了数学学科区别于其他学科的核心素养，课程目标也从“双基”到“四基”、“两能”到“四能”。课程标准不仅注重培养学生的数学能力和素质，而且对中小学数学教师提出了更高的要求，因此培养师范生专业素质和能力也成为了学校教师队伍建设的重心所在.

近日，教育部印发《中学教育专业师范生教师职业能力标准（试行）》等五个文件。该文件明确了师范生的基本素质，为进一步加强师范院校的学科建设、提升师范生教学能力和专业水平提供了依据。而数学教学技能是数学师范生教育教学能力水平的重要体现，因此，了解目前师范生的技能掌握状况，对师范院校的数学教育专业发展具有重要意义。



\section{研究内容}

本论文研究的主要内容有：（1）分析数学专业师范生教学技能掌握现状；（2）揭示数学专业师范生在教学技能方面的不足，从理论上和实践上为高师培训计划的修改和课程设置提供参考；（3）提出促进数学专业师范生教学技能培养的合理建议。

首先从内涵和外延定义相关概念，确定研究的出发点和立足点；特别是对教学技能的定义，对教学技能进行细致全面的分类，使研究精细化、专业化和科学化；了解现有的研究状况，总结、归纳和分析与本论文有关的研究成果，确定研究展开的起点；在此基础上设计调查问卷；找出数学专业师范生教学技能发展的不足之处，在此基础上有针对性地寻求解决方案，提供相关建议。



\section{研究方法}

本研究主要采用文献分析法、调查问卷和访谈法.

文献分析法：检索相关文献，查阅相关论文，整理分析国内外的研究成果，充分研究已有的研究成果，明确相关问题研究的现状，发现目前需要研究的具体问题。文献研究主要包括文字资料的搜集、整理、处理、筛选、归纳。

调查问卷和访谈法：首先对数学专业师范生、中学数学教师、数学教育专家进行访谈，在此基础上编制调查问卷。以此问卷作为研究工具，对华中师范大学2018级和2019级数学专业师范生进行调查研究，为研究收集有价值的数据。



\section{研究意义}

数学教师最重要的职业素养是具备良好的数学教学技能，它关系到教学质量的提高、教师专业的发展和学生的数学学习。本文的研究结果能够反映出师范生的数学教学技能掌握状况，并深入探究数学教学技能在年龄、性别等方面的差异，了解其内部各维度的相关性，从理论上和实践上为高师培训计划的修改和课程设置提供参考，从而促进师范生数学教学技能的提高。